\section*{Capítulo \ref{ch:g12:deriv} --- Derivación y convexidad}

\begin{solution}{exo:g12:deriv:1}
$f'(x) = 3x^2 - 5$.

Regla del cociente:
$g'(x) = \dfrac{(x^2+1) - x\cdot 2x}{(x^2+1)^2} = \dfrac{1 - x^2}{(x^2+1)^2}$.

Regla de la cadena con $u = 2x+1$:
$h'(x) = 5 \cdot 2 \cdot (2x+1)^4 = 10(2x+1)^4$.

Regla de la cadena con $u = x^2 + 1$:
$k'(x) = \dfrac{2x}{2\sqrt{x^2+1}} = \dfrac{x}{\sqrt{x^2+1}}$.
\end{solution}

\begin{solution}{exo:g12:deriv:2}
$f'(x) = 2x - 3$, luego $f(2) = -1$ y $f'(2) = 1$: la \omterm{def:g12:deriv:tangent}{tangente} es
$y = -1 + (x - 2) = x - 3$. La separación es
\[
f(x) - (x - 3) = x^2 - 4x + 4 = (x-2)^2 \geq 0 ,
\]
así que la curva queda por encima de la \omterm{def:g12:deriv:tangent}{tangente} en todas partes y solo la
toca en $x = 2$ (como cabía esperar: $f$ es \omterm{def:g12:deriv:convex}{convexa}).
\end{solution}

\begin{solution}{exo:g12:deriv:3}
Los dos son cocientes incrementales. Con $f(x) = (1+x)^{100}$ en $0$:
$f'(x) = 100(1+x)^{99}$, así que el límite es $f'(0) = 100$. Con
$g(x) = \sqrt{x}$ en $4$: $g'(x) = \frac{1}{2\sqrt{x}}$, así que el límite
es $g'(4) = \frac14$.
\end{solution}

\begin{solution}{exo:g12:deriv:4}
\omterm{def:g11:func:function}{Dominio} $\R \setminus \{1\}$. Límites: en $\pm\infty$,
$f(x) \sim x \to \pm\infty$; en $1^\pm$, el numerador $\to 4 > 0$ y el
denominador $\to 0^\pm$, luego $f \to \pm\infty$; la recta $x = 1$ es
\omterm{def:g12:limcont:asymptote}{asíntota vertical}.

\[
f'(x) = \frac{2x(x-1) - (x^2+3)}{(x-1)^2} = \frac{x^2 - 2x - 3}{(x-1)^2}
= \frac{(x-3)(x+1)}{(x-1)^2}.
\]
Por tanto, $f' > 0$ en $\intoo{-\infty}{-1}$ y en $\intoo{3}{+\infty}$, y
$f' < 0$ en $\intoo{-1}{1}$ y en $\intoo{1}{3}$. La \omterm{def:g11:func:function}{función} crece hasta un
\omterm{def:g12:deriv:extremum}{máximo relativo} $f(-1) = -2$, decrece hasta $-\infty$, salta a $+\infty$
tras la \omterm{def:g12:limcont:asymptote}{asíntota}, decrece hasta un \omterm{def:g10:functions:extrema}{mínimo} relativo $f(3) = 6$ y crece
después.
\end{solution}

\begin{solution}{exo:g12:deriv:5}
Para $x \in \intoo{0}{15}$, la caja tiene una base cuadrada de lado
$30 - 2x$ y altura $x$, así que
\[
V(x) = x(30 - 2x)^2 .
\]
$V'(x) = (30-2x)^2 + x \cdot 2(30-2x)(-2) = (30-2x)(30 - 2x - 4x)
= (30-2x)(30-6x)$. En $\intoo{0}{15}$, $30 - 2x > 0$, así que $V'$ tiene el
signo de $30 - 6x$: positivo antes de $x = 5$ y negativo después. El
volumen es \omterm{def:g10:functions:extrema}{máximo} para $x = 5$ cm, y vale
$V(5) = 5 \times 20^2 = 2000\ \text{cm}^3$.
\end{solution}

\begin{solution}{exo:g12:deriv:6}
\emph{1.} $f'(x) = 4x^3 - 12x^2$ y $f''(x) = 12x^2 - 24x = 12x(x-2)$. Así
que $f'' > 0$ en $\intoo{-\infty}{0}$ y en $\intoo{2}{+\infty}$ (\omterm{def:g12:deriv:convex}{convexa}),
y $f'' < 0$ en $\intoo{0}{2}$ (\omterm{def:g12:deriv:convex}{cóncava}). $f''$ cambia de signo en $0$ y en
$2$: los dos son \omterm{def:g12:deriv:inflection}{puntos de inflexión}.

\emph{2.} En $a = 2$: $f(2) = 16 - 32 + 10 = -6$, $f'(2) = 32 - 48 = -16$ y
la \omterm{def:g12:deriv:tangent}{tangente} es $y = -6 - 16(x-2) = -16x + 26$. La separación es
\[
g(x) = f(x) + 16x - 26 = x^4 - 4x^3 + 16x - 16 .
\]
Como $g(2) = g'(2) = g''(2) = 0$, $(x-2)^3$ divide a $g$; la división da
$g(x) = (x-2)^3(x+2)$. Cerca de $x = 2$, $x + 2 > 0$, así que $g$ tiene el
signo de $(x-2)^3$: negativo antes de $2$ y positivo después. La curva
atraviesa la \omterm{def:g12:deriv:tangent}{tangente}: $2$ es, en efecto, un \omterm{def:g12:deriv:inflection}{punto de inflexión}.
\end{solution}

\begin{solution}{exo:g12:deriv:7}
La \omterm{def:g11:func:function}{función} $f(x) = \sqrt{x}$ cumple $f''(x) = -\frac{1}{4}x^{-3/2} < 0$ en
$\intoo{0}{+\infty}$: es \omterm{def:g12:deriv:convex}{cóncava}, así que su \omterm{def:g10:functions:graph}{gráfica} queda \emph{por
debajo} de todas sus \omterm{def:g12:deriv:tangent}{tangentes}. La \omterm{def:g12:deriv:tangent}{tangente} en $a = 1$ es
\[
y = f(1) + f'(1)(x - 1) = 1 + \tfrac12 (x-1) = \frac{x+1}{2},
\]
de donde $\sqrt{x} \leq \frac{x+1}{2}$ para todo $x > 0$. La igualdad
significa que la curva toca la \omterm{def:g12:deriv:tangent}{tangente}, cosa que solo ocurre en el punto
de tangencia $x = 1$. (Equivalentemente:
$\left(\sqrt x - 1\right)^2 \geq 0$.)
\end{solution}

\begin{solution}{exo:g12:deriv:8}
Por el \cref{thm:g12:deriv:convexity}, la \omterm{def:g10:functions:graph}{gráfica} de una \omterm{def:g11:func:function}{función} \omterm{def:g12:deriv:convex}{convexa}
\omterm{def:g12:deriv:derivative}{derivable} queda por encima de todas sus \omterm{def:g12:deriv:tangent}{tangentes}. La \omterm{def:g12:deriv:tangent}{tangente} en $a$ es
horizontal ($f'(a) = 0$), de \omterm{def:g10:algebra:equation}{ecuación} $y = f(a)$. Por tanto,
$f(x) \geq f(a)$ para todo $x \in \R$: $f(a)$ es el \omterm{def:g10:functions:extrema}{mínimo} absoluto.
\end{solution}

\begin{solution}{exo:g12:deriv:9}
\emph{Perímetro fijo.} Un rectángulo de lados $x$ y $\frac{p}{2} - x$
($0 < x < \frac p2$) tiene área $S(x) = x\left(\frac p2 - x\right)$.
Entonces $S'(x) = \frac p2 - 2x$ se anula en $x = \frac p4$, es positiva
antes y negativa después: el \omterm{def:g10:functions:extrema}{máximo} está en $x = \frac{p}{4}$, donde los
dos lados valen $\frac p4$; un cuadrado.

\emph{Área fija.} Los lados $x$ y $\frac Ax$ dan un perímetro
$P(x) = 2\left(x + \frac Ax\right)$, con
$P'(x) = 2\left(1 - \frac{A}{x^2}\right)$, negativa para $x < \sqrt A$ y
positiva después: \omterm{def:g10:functions:extrema}{mínimo} en $x = \sqrt{A}$; otra vez un cuadrado.

\emph{Relación.} Las dos afirmaciones son duales. Supongamos que un
rectángulo no cuadrado minimizara el perímetro con área fija $A$; el
cuadrado de ese mismo perímetro tendría área $> A$ por la primera
afirmación, y encogerlo por semejanza hasta el área $A$ reduciría
estrictamente su perímetro, en contradicción con la minimalidad. Cada
afirmación implica, pues, la otra.
\end{solution}

\begin{solution}{pb:g12:deriv:1}
\textbf{1.} $30x\left(3x^2 + 1\right)^4$; \quad
$\dfrac{x}{\sqrt{x^2 + 9}}$; \quad
$-\dfrac{2x}{\left(x^2 + 1\right)^2}$.

\textbf{2.} $y(4) = 5$, \omterm{def:g10:reffunc:affine}{pendiente} $\frac45$: \omterm{def:g12:deriv:tangent}{tangente}
$y = 5 + \frac45(x - 4)$, es decir, $y = \frac45 x + \frac95$.

\textbf{3.} $f'(x) = \frac{(x^2 + 1) - x \cdot 2x}{(x^2+1)^2}
= \frac{1 - x^2}{(x^2 + 1)^2}$: negativa, positiva y negativa al pasar por
$-1$ y $1$: \omterm{def:g10:functions:extrema}{mínimo} $f(-1) = -\frac12$, \omterm{def:g10:functions:extrema}{máximo} $f(1) = \frac12$, con
\omterm{def:g12:limcont:asymptote}{asíntota horizontal} $0$ a los dos lados.

\textbf{4.} $g'(x) = 3x^2 - 6x = 3x(x - 2)$: \omterm{def:g12:seq:monotonic}{creciente}, \omterm{def:g10:functions:variations}{decreciente} y
\omterm{def:g12:seq:monotonic}{creciente} al pasar por $0$ y $2$; \omterm{def:g12:deriv:extremum}{máximo relativo} $g(0) = 4$ y \omterm{def:g10:functions:extrema}{mínimo}
relativo $g(2) = 0$. $g''(x) = 6x - 6$: \omterm{def:g12:deriv:convex}{cóncava} antes de $1$ y \omterm{def:g12:deriv:convex}{convexa}
después: inflexión en $(1, 2)$.

\textbf{5.} \omterm{def:g12:deriv:tangent}{Tangente} en $1$: \omterm{def:g10:reffunc:affine}{pendiente} $g'(1) = -3$:
$y = 2 - 3(x - 1) = -3x + 5$. \omterm{def:g11:seq:arithmetic}{Diferencia}:
$x^3 - 3x^2 + 4 - (-3x + 5) = x^3 - 3x^2 + 3x - 1 = (x - 1)^3$, que cambia
de signo en $1$: la \omterm{def:g12:deriv:tangent}{tangente} \emph{atraviesa} la curva; ese es el
comportamiento característico en un \omterm{def:g12:deriv:inflection}{punto de inflexión}, donde la curva
cambia de lado.

\textbf{6.} $T(x) = \dfrac{\sqrt{x^2 + 900}}{5} +
\dfrac{\sqrt{(40 - x)^2 + 400}}{2}$. Entrar antes de $0$ o más allá de
$40$ alarga \emph{los dos} tramos: no tiene sentido.

\textbf{7.} $T'(x) = \dfrac{x}{5\sqrt{x^2 + 900}} -
\dfrac{40 - x}{2\sqrt{(40 - x)^2 + 400}}$.

\textbf{8.} En el triángulo del tramo a la carrera, el lado paralelo a la
orilla mide $x$ y la hipotenusa, $\sqrt{x^2 + 900}$: el \omterm{def:g11:trigo:cossin}{seno} del ángulo con
la perpendicular es precisamente su cociente (cateto opuesto entre
hipotenusa); y del mismo modo,
$\sin\theta_2 = \frac{40 - x}{\sqrt{(40-x)^2 + 400}}$. Entonces
$T'(x) = 0$ se lee $\frac{\sin\theta_1}{5} = \frac{\sin\theta_2}{2}$: la
ley de Snell con velocidades $5$ y $2$.

\textbf{9.} $T(24) \approx 20.5$ s; $T(40) = 10 + 10 = 20$ s;
$T(0) = 6 + \sqrt{2000}/2 \approx 28.4$ s. La línea recta pierde incluso
frente a ``correr todo el camino y después nadar en línea recta'', y las
dos pierden frente al camino refractado.

\textbf{10.} La \omterm{thm:g12:limcont:ivt}{dicotomía} sobre $T'$ (negativa en $24$ y positiva en $40$)
\omterm{def:g12:seq:limit}{converge} a $x^* \approx 34$ m, con $T(x^*) \approx 19.5$ s: un segundo más
rápido que la línea recta; la \omterm{def:g11:seq:arithmetic}{diferencia} entre un rescate y una tragedia.

\textbf{11.} El tiempo de cada tramo es una \omterm{def:g11:func:function}{función} \omterm{def:g12:deriv:convex}{convexa} de $x$ (una
rama de la \omterm{def:g11:quad:discriminant}{raíz} cuadrada de un polinomio de segundo grado), así que $T$ es
\omterm{def:g12:deriv:convex}{convexa},
y un punto crítico de una \omterm{def:g11:func:function}{función} \omterm{def:g12:deriv:convex}{convexa} \omterm{def:g12:deriv:derivative}{derivable} es su \omterm{def:g10:functions:extrema}{mínimo} absoluto
(\cref{exo:g12:deriv:8}): $x^*$ no es solo estacionario, es imbatible.

\textbf{12.} La luz es más lenta en el agua; por el principio de Fermat, el
camino más rápido del aire al agua se quiebra en la superficie exactamente
según la ley de Snell, así que los rayos de un lápiz sumergido llegan al
ojo quebrados y el cerebro, que prolonga líneas rectas, ve el lápiz roto en
la línea del agua.

\textbf{13.} $(x^2)'' = 2 > 0$: \omterm{def:g12:deriv:convex}{convexa}. \omterm{def:g12:deriv:convex}{Convexidad} en el \omterm{prop:g10:coordgeom:midpoint}{punto medio}:
$f\left(\frac{a+b}{2}\right) \leq \frac{f(a) + f(b)}{2}$, que es lo que se
afirma. A mano:
$\frac{a^2 + b^2}{2} - \left(\frac{a+b}{2}\right)^2 =
\frac{(a - b)^2}{4} \geq 0$.

\textbf{14.} Tomando raíces cuadradas (que es \omterm{def:g12:seq:monotonic}{creciente}) en la pregunta 13:
$\frac{a + b}{2} \leq \sqrt{\frac{a^2 + b^2}{2}}$. Cadena completa con $2$
y $8$: armónica $\frac{2 \cdot 16}{10} = 3.2$; geométrica $\sqrt{16} = 4$;
aritmética $5$; cuadrática $\sqrt{34} \approx 5.83$: cada \omterm{def:g11:stat:mean}{media} se inclina
ante la siguiente.

\textbf{15.} \omterm{def:g12:deriv:tangent}{Tangente} a $\frac1x$ en $1$: valor $1$ y \omterm{def:g10:reffunc:affine}{pendiente} $-1$:
$y = 2 - x$. Las curvas \omterm{def:g12:deriv:convex}{convexas} se sitúan por encima de sus \omterm{def:g12:deriv:tangent}{tangentes}
(\cref{thm:g12:deriv:convexity}): $\frac1x \geq 2 - x$ para todo $x > 0$,
con igualdad exactamente en el punto de contacto $x = 1$.

\textbf{16.} En la inflexión, el recuento \emph{diario} de casos (la
\omterm{def:g12:deriv:derivative}{derivada}) alcanza su \omterm{def:g10:functions:extrema}{máximo}: el crecimiento deja de acelerarse y empieza a
frenar; es la primera señal matemática de que la ola está girando, mucho
antes de que caigan los propios recuentos. Antes de ella, $f'' > 0$ (cada
día peor que el anterior); después, $f'' < 0$ (sigue creciendo, pero más
despacio).

\textbf{17.} $S'(r) = 4\pi r - \frac{2V}{r^2} = 0$ da
$r^3 = \frac{V}{2\pi} = \frac{330}{2\pi}$: $r \approx 3.74$ cm, y
$h = \frac{330}{\pi r^2} \approx 7.49$ cm.

\textbf{18.} $h = \frac{V}{\pi r^2} = \frac{2\pi r^3}{\pi r^2} = 2r$ en el
óptimo: la altura es igual al diámetro, sea cual sea el volumen.

\textbf{19.} $S''(r) = 4\pi + \frac{4V}{r^3} > 0$: \omterm{def:g12:deriv:convex}{convexa}, así que el
radio crítico da el \omterm{def:g10:functions:extrema}{mínimo} absoluto. Las latas reales son más altas por el
agarre, el apilado, la presencia en el lineal y el metal más grueso de la
tapa y el fondo: la optimización siempre minimiza exactamente lo que has
escrito, no lo que querías decir.

\textbf{20.} Socorrista: modelizar $T(x)$, derivar, Snell, \omterm{def:g12:deriv:convex}{convexidad}, un
segundo ganado. Lata: modelizar $S(r)$, derivar, $h = 2r$, \omterm{def:g12:deriv:convex}{convexidad},
metal ahorrado. Viga (\cref{pb:g11:deriv:1}): modelizar $S(w)$, derivar,
$d = w\sqrt2$, tabla, rigidez ganada. Y, por el principio de Fermat, la
propia luz recorre esta cadena en cada superficie que encuentra: la
naturaleza fue el primer optimizador.
\end{solution}
